En el caso sin ruido un retraso generalmente requiere para la
codificacion\'on ideal ahora esta tiene una funci\'on adicional que
permite una muestra larga de ruido para afectar la se\~{n}al antes de
un juicio, se hace en el punto de resepci\'on para el mensaje
original. Incrementando el tama\~{n}o del ejemplo siempre agudiza las
posibles afirmaciones estad\'isticas.

El contenido del teorema \ref{th:11} y su prueba se pueden formular de
una manera diferente que muestra una conexi\'on sin ruido de una
manera mas clara. Considere las posibles duraciones de las se\~{n}ales
T y supongamos un subconjutno de ellos se seleccionan para ser
usados. Los que esten en el subconjunto se utilizen todos con igual
probabilidad, y suponiendo que el resector es construido para
seleccionar, como la se\~{n}al original. la causa m\'as probable de
subconjutno cuando una se\~{n}al perturbada es recivida. Nosotros
definimos $N(T,q)$ siendo el numero m\'aximo de se\~{n}ales que
podemos elegir para el subconjunto tal que la probabilidad de una
interpretaci\'{o}n incorrecta sea menor o igual a $q$.

\begin{theorem}
\label{th:12}
  \begin{equation} \lim_{T \to{} \infty}\frac{\log{N}(T,q)}{T} =
    C, \end{equation} donde $C$ es la capacidad de canal,a condici\'on
    que q no sea igual a 0 o 1
\end{theorem}

En otras palabras, no importa lo que nos propusimos de limites de
confiabilidad, podemos distinguir de forma fiable en el tiempo $T$
suficientes mensajes para corresponder a $CT$ bits, cuando $T$ es
suficientemente grande. En el teorema \ref{t12} podemos comparar la
capacidad de un canal sin ruido en la secci\'on uno PONER REF.

\clearpage

\chapter{Ejemplo de un canal discreto y su capacidad}
\label{sec:15}

Un ejemplo simple de un canal discreto se indica en la
figura \ref{fig:11}. Hay tres posibles s\'imbolos. El primero nunca se
vera afectado por el ruido. El segundo y la tercera tiene cada una
probabilidad $p$ de que viene a trav\'es impertubable q de ser
cambiado en el otro elemento par. Nosotros tenemos (dejar $\alpha =
-[plogp + q\log{q}]$ y $p$ y $q$ son probabilidades de estar usando
los s\'imbolos primero y segundo).

\begin{figure}[!ht]
\centerline{\includegraphics[width=80mm]{ejemplo.png}}
\caption{Un ejemplo de un canal discreto.}
\label{fig:11}
\end{figure}

\begin{equation}
\begin{array}{rcl}
H(x) &=& -P\log{P} - 2Q\log{Q} \\
H_y(x) &=& 2Q\alpha
\end{array}
\end{equation}

Nosotros deseamos elegir $P$ y $Q$, de tal manera que se maximice
$H(x) - H_y(x)$ sujeto a la restricci\'on $P + 2Q = 1$.

Por lo tanto consideremos:
\begingroup
\renewcommand*{\arraystretch}{2.0}
\begin{equation}
\begin{array}{rcl}
U &=& -P\log{P} - 2Q\log{Q} -2Q\alpha + \lambda(P+2Q) \\
\displaystyle\frac{{\partial U}}{{\partial P}} &=& -1 - \log{P} + \lambda = 0 \\
\displaystyle\frac{{\partial U}}{{\partial P}} &=& -2 - 2\log{Q} -2\alpha + 2\lambda = 0
\end{array}
\end{equation}
\endgroup

Eliminando $\lambda$
\begin{equation}
\begin{array}{rcl}
\log{P} &=& \log{Q} + \alpha \\
P &=& Q e^\alpha = Q\beta 
\end{array}
\end{equation}

\begin{equation}
  P = \frac{\beta}{\beta + 2}   Q = \frac{1}{\beta + 2}
\end{equation}

La capacidad del canal es de:

\begin{equation}
  C = \log{\frac{\beta + 2}{\beta}}
\end{equation}

Notese como esto comprueba los valores obvios en los casos: $p = 1$ y $p = \frac{1}{2}$. En primero, $\beta = 1$ y $C = \log{3}$, 
el cual es correcto debido a que el canal es entonces sin ruido con tres posibles s\'imbolos. Si $p = \frac{1}{2}$, $\beta = 2$ y 
$C = \log{2}$.Aqu\'i el segundo y el tercer s\'imbolos, no se pueden distinguir en absoluto y act\'uan conjuntamente como un solo 
s\'imbolo. El primer s\'imbolo se utiliza con una probabilidad $P = \frac{1}{2}$  y el segundo junto al tercero con probabilidad $\frac{1}{2}$.
Esto puede ser distribuido entre ellos de cualquier modo deseado y todav\'ia alcanzar la m\'axima capacidad.

Para los valores intermedios de la capacidad del canal $p$ estara entre $\log{2}$ y $\log{3}$.Esta distinci\'on 
entre el segundo y tercer s\'imbolo transmite alguna informaci\'on, pero no tanto como en el caso sin ruido.
El primer s\'imbolo se utiliza tanto mas frecuentemente que los otros dos, debido a su ausencia de ruido.

\clearpage

\chapter{La capacidad de canales en cietos casos especiales}
\label{sec:16}

Si el ruido afecta s\'imbolos sucesivos de canal de forma
independiente pueden ser descritos por un conjunto de transici\'on de
probabilidades $p_{i,j}$.Esta es la probabilidad, si el simbolo i es
enviado, que j sera recibido.La tasa de canal m\'aximo viene dado por
el m\'aximo de:
\begin{equation}
  \sum_{i,j}P_i p_{i,j} \log{\sum{P_i p_{i,j}}} + \sum_{i,j}P_i p_{i,j}\log{p_{i,j}}
\end{equation}
 
Multiplicar por $P_s$ y sumando en $s$ muestra que $\mu = C$. Vamos a la inversa de $p_{sj}$ (si existe) en $h_{st}$  de modo que 
$\sum_{s}$ $h_{st}$ $p_{sj} = \delta_{tj}$. Entonces: 
\begin{equation}
  \sum_{s,j}h_{st} p_{s,j} \log{p_{s.j}} - \log{\sum_{i}P_i p_{i,t}} = C \sum_{s} h_{s,t}.
\end{equation}
Por lo tanto:
\begin{equation}
  \sum_{i} P_i p_{i,t} = exp[- C \sum_{s} h_{s,t}+ \sum_{s,j} h_{s,t} p_{s,j} \log{p_{s,j}}]
\end{equation}
o  
\begin{equation}
  P_i = \sum_{t} h_{i,t} exp[ - C \sum_{s} h_{s,t}+ \sum_{s,j} h_{s,t} p_{s,j} \log{p_{s,j}} ].
\end{equation}
Este es el sistema de ecuaciones para determinar el valor m\'aximo de
$P$, con $C$ se determina de manera que $\sum P_i = 1$. Cuando esto
esta hecho, $C$ sera la capacidad del canal y $P_i$ las probabilidades
para los s\'imbolos de canal para lograr esta capacidado.  Si cada
s\'imbolo de entrada tiene el mismo conjunto de probabilidades en las
l\'ineas que emergen de ella y lo mismo sucede a cada s\'imbolo de
salida, la capacidad puede ser calculada f\'acilmente. Los ejemplos se
muestran en la figura 12.  En tal caso $H_xy$ es independiente de la
destribuci\'on de probabilidades de los s\'imbolos de entrada, y esta
dad por $-\sum p_i \log{p_i}$.  Cuando $p_i$ son los valores de
probabilidad de transici\'on de cualquier s\'imbolo de entrada.La
capacidad del canal es
\begin{equation}
  \max [H(y) - H_x(y)] = \max H(y) + \sum p_i \log{p_i}.
\end{equation}
El valor m\'aximo de $H(y)$ esta claramente  $\log{m}$ donde $m$ es el numero de s\'imbolos de salida, ya que es posible para hacer 
todos igualmente probables haciendo los s\'imbolos de entradas igualmente probables. La capacidad del canal es por lo tanto
\begin{equation}
  C = \log{m} + \sum p_i \log{p_i}.
\end{equation}

\begin{figure}
  \centering
  \begin{subfigure}[b]{0.4\textwidth}
    \centering
    \includegraphics[width=\textwidth]{ejemplo.png}
    \caption{Primer ejemplo.}
    \label{fig12a}
  \end{subfigure}%
\quad
  \centering
  \begin{subfigure}[b]{0.4\textwidth}
    \centering
    \includegraphics[width=\textwidth]{ejemplo.png}
    \caption{Segundo ejemplo.}
    \label{fig12b}
  \end{subfigure}%
\quad
  \centering
  \begin{subfigure}[b]{0.4\textwidth}
    \centering
    \includegraphics[width=\textwidth]{ejemplo.png}
    \caption{Tercer ejemplo.}
    \label{fig12c}
  \end{subfigure}%
\caption{Ejemplos de canales discretos con las mismas probabilidades de
 transici\'{o}n para cada entrada y cada salida.}
\label{fig:12}
\end{figure}

En la figura \ref{fig12a} ser\'ia
\begin{equation}
  C = = \log{4} - \log{2} 0 \log{2}.
\end{equation}
Esto se podr\'ia lograr mediante el uso solo el primer y el tercer
s\'{\i}mbolo. En la Figura \ref{fig12b}:
\begin{equation}
  C = \log{4} - \frac{2}{3}\log{3} - \frac{1}{3}\log{6}
  = \log{4} - \log{3} - \frac{1}{3}\log{2}
  = \log{\frac{1}{3}} 2^{\frac{5}{3}}
\end{equation}
En la figura \ref{fig12c} nosotros tenemos:
\begin{equation}
  C = \log{3} - \frac{1}{2}\log{2} - \frac{1}{3}\log{3} - \frac{1}{6} \log{6}
  = \log  {\frac{3}{ 2^{\frac{1}{2}}}} 3^{\frac{1}{3} 6^{\frac{1}{6} }}.
\end{equation}
Supongamos que los s\'imbolos se dividen en varios grupos tal que el
ruido causa a un s\'imbolo en un grupo a ser confundido con un
s\'imbolo de otro grupo. Deja la capacidad de un grupo n-\'esimo ser
$C_n$ (en bits por segundo) donde solo utilizamos los s\'imbolos de
este grupo. Entonces es f\'acil desmotrar de todo el conjunto, la
probabilidad $P_n$ para todos los s\'imbolos del grupo $n$-\'{e}esimo
deber\'ia ser:
\begin{equation}
  P_n = \frac{2^{C_n}}{\sum 2^{C_n}}
\end{equation}

En un grupo la probabilidad se distribuye tal como ser\'ia si estos eran los \'unicos s\'imbolos que se utilizan. 
La capacidad del canal es:
\begin{equation}
  C = \log{\sum 2^{C_n}}.
\end{equation}

\clearpage

\chapter{Un ejemplo de codificaci\'on eficiente}
\label{sec:17}

El siguiente ejemplo, aunque es un poco realista, es un caso en que la
coincidencia exacta para un canal con ruido, es posible. Hay dos
s\'imbolos de canal $0$ y $1$, y el ruido les afecta en bloques de siete
s\'imbolos.  Un bloque de siete o se transmite sin error, o
exactamente un s\'imbolo de los siete es incorrecta.Estas ocho
posibilidades son igualmente probables. Nosotros tenemos:
\begin{equation}
C = \max [ H(y) - H_x(y) ]
= \frac{1}{7}[7 + \frac{8}{8}\log{\frac{1}{8}}]
= \frac{4}{7} \text{ bits}/\text{s\'{\i}mbolos}
\end{equation}
Un c\'odigo eficiente, permite la correci\'on de todos los errores y
transmitir a la tasa $C$, es el siguiente (encontrado por un m\'etodo
de R.\ Hamming):
